\chapter{Linux Tuning}\label{ch:ll:linux-tuning}

Run a thread that does nothing but read the \omterm{def:ll:cpu-microarchitecture:tsc}{time-stamp counter} in a loop, and record every gap between two consecutive
readings longer than \qty{200}{\nano\second}, when one pass of the loop takes about fifteen. Each such gap is a moment when
the thread was not running: something else had the CPU. On the laptop of this book, pinned to a CPU that no other program
of ours uses, the thread loses about 1.3\% of its time: some 750 times a second for five microseconds or more, a few
times a second for more than a hundred, and every few seconds for a millisecond or more. A market-data
message that arrives during a gap waits for the rest of it. None of this is the program's fault, and none of it can be
fixed in the program: it is the operating system doing its own work on the CPU the program thought it had. This chapter is
about taking the CPU back. It moves the kernel's work, the interrupts and the timer tick off the cores that run the hot
path, replaces sleeping with polling, keeps the program's memory in RAM, and ends with a build that audits a host against
the placement plan of chapter~4 before a trading process is allowed to start on it.

\section{Core isolation and CPU affinity}

\begin{definition}[CPU affinity, core isolation, housekeeping core]\label{def:ll:linux-tuning:isolation}
The \emph{CPU affinity}\index{CPU affinity} of a thread is the set of CPUs the scheduler may run it on; it is set with
\texttt{sched\_setaffinity} (or \texttt{taskset} from a shell) and inherited by the threads a thread creates. \emph{Core
isolation}\index{core isolation} removes CPUs from the scheduler's load balancing, so that no thread runs on them unless its
affinity names them explicitly, and moves the kernel's own unbound work away from them. A \emph{housekeeping
core}\index{housekeeping core} is a CPU left out of the isolated set, which runs everything else: the kernel's work, the
interrupts of devices the hot path does not use, and the process's cold threads.
\end{definition}

Affinity alone fixes where a thread runs, not what else runs there: a pinned thread still shares its CPU with any other
runnable task the scheduler places on it, with kernel threads, and with interrupts. The kernel documentation calls all of
these ``noise'', asynchronous (interrupts, timers, work queues, kernel threads) or synchronous (system calls and \omterm{def:ll:cpp-for-latency-i-memory:fault}{page
faults}), and notes that it usually goes unnoticed: the timer interrupt can run 1\,024 times a second ``without a significant
and measurable impact most of the time''. The exceptions it names are workloads such as very low latency network processing.
Isolation is how the noise is moved: the boot parameter \texttt{isolcpus} (with its default \texttt{domain} flag) takes CPUs
out of ``the general SMP balancing and scheduling algorithms'', irreversibly until the next boot, and control-group cpusets
do the same at run time. A thread reaches an isolated CPU only through its affinity, which is exactly what chapter~9's
\texttt{firm.affinity} sets from chapter~4's plan.

\Cref{fig:ll:linux-tuning:hiccups} shows what the hook's thread sees in three situations on this untuned laptop. Pinned
alone or left free to move, it loses about one percent of its time in gaps of microseconds, and a few gaps a second reach
a hundred microseconds or more. Sharing its CPU with a second spinning thread, it loses about half of it, about 120 times a
second for about four milliseconds each: the scheduler hands the CPU to the rival for a slice, and a slice of that length
matches the period of this kernel's tick, which runs 250 times a second (its configuration, \texttt{CONFIG\_HZ}). The shared case is the one that
isolation removes outright; the small gaps of the other two need the rest of this chapter, and some of them, on a virtual
machine, belong to the hypervisor, which the guest cannot see or tune.

\begin{omfigure}
\begin{tikzpicture}
\begin{axis}[width=11.5cm,height=5.8cm,xmode=log,ymode=log,log origin y=infty,xmin=150,xmax=1.3e7,ymin=0.1,ymax=20000,
  xlabel={gap length $x$ (ns, log scale)},ylabel={gaps per second $\ge x$ (log scale)},
  tick label style={font=\small},label style={font=\small},grid=major,grid style={gray!25},
  unbounded coords=discard,filter discard warning=false,
  xtick={1e3,1e4,1e5,1e6,1e7},xticklabels={\qty{1}{\micro\second},\qty{10}{\micro\second},\qty{100}{\micro\second},
  \qty{1}{\milli\second},\qty{10}{\milli\second}},
  legend image code/.code={\draw[#1] (0cm,0cm) -- (0.6cm,0cm);},
  legend columns=3,legend style={font=\footnotesize,at={(0.5,-0.3)},anchor=north,draw=none,fill=none,
  /tikz/every even column/.append style={column sep=8pt}},table/col sep=comma]
\addplot[omDef,very thick,mark=*,mark size=1.2pt] table[x=threshold_ns,y=alone]{figdata/low-latency/13-linux-tuning/measured_hiccups.csv};
\addplot[omMeth,thick,dashed,mark=square*,mark size=1.2pt] table[x=threshold_ns,y=unpinned]{figdata/low-latency/13-linux-tuning/measured_hiccups.csv};
\addplot[omThm,very thick,mark=triangle*,mark size=1.6pt] table[x=threshold_ns,y=shared]{figdata/low-latency/13-linux-tuning/measured_hiccups.csv};
\legend{pinned alone,unpinned,{pinned, sharing its CPU}}
\end{axis}
\end{tikzpicture}

\omcaption{The hiccup meter: how often a thread spinning on the \omterm{def:ll:cpu-microarchitecture:tsc}{time-stamp counter} is kept from running for at least $x$,
over five seconds, pinned to a CPU no other program of ours uses, free to move, and pinned to a CPU shared with a second
spinning thread. Measured on a laptop (Intel Core Ultra 7 155H) under WSL2, no isolated cores, benchmark at normal
priority on an otherwise idle machine. Data: \texttt{bench\_tuning.py}.}\label{fig:ll:linux-tuning:hiccups}
\end{omfigure}

On a tuned host the plan of chapter~4 becomes a boot configuration. \Cref{fig:ll:linux-tuning:host} draws it for the
two-socket fixture of that chapter: the three hot threads on three cores of the network card's socket, their hyper-threading
siblings taken offline (the kernel's own checklist says to avoid SMT, ``to prevent your hardware thread from being
`preempted' by another one''), the card's interrupts on housekeeping CPUs of the same socket, and everything else on the
rest. The documentation asks for at least one housekeeping CPU, preferably one per \omterm{def:ll:multi-socket-machines-and-interconnects:numa}{NUMA node}; a trading host keeps several,
since the logger, the monitoring agents and the kernel all live there, and a housekeeping CPU that is itself overloaded delays
the work (log draining, timers) that the hot path hands to it.

\begin{omfigure}
\begin{tikzpicture}[font=\scriptsize,cpu/.style={draw=gray!70,minimum width=0.62cm,minimum height=0.5cm,inner sep=0pt}]
\foreach \s/\y in {0/0,1/-1.9} {
  \node[anchor=east,font=\footnotesize] at (-0.2,\y-0.3) {socket \s};
  \foreach \k in {0,...,7} {
    \pgfmathtruncatemacro{\a}{16*\s+\k}\pgfmathtruncatemacro{\b}{16*\s+\k+8}
    \node[cpu,fill=gray!12] (c\a) at (0.7*\k+0.3,\y) {\a};
    \node[cpu,fill=gray!12] (c\b) at (0.7*\k+0.3,\y-0.6) {\b};
  }
}
\foreach \c in {16,17,18,19,20} {\node[cpu,fill=omDef!30] at (c\c) {\c};}
\foreach \c in {21,22,23} {\node[cpu,fill=omThm!55,text=white,font=\scriptsize\bfseries] at (c\c) {\c};}
\foreach \c in {29,30,31} {\node[cpu,fill=white,draw=gray!70,dashed,text=gray] at (c\c) {\c};}
\node[draw=omDef,thick,fill=omDef!8,align=center] (nic) at (1.7,-3.55) {network card, node 1};
\draw[omDef,thick,-Stealth] (nic.north) -- node[right,font=\scriptsize]{interrupts to CPUs 16--20} (1.7,-2.75);
\foreach \i/\f/\d/\t in {0/omThm!55/solid/{isolated, tickless: feed, strategy, gateway},1/white/dashed/{their siblings, offline},
  2/omDef!30/solid/{the card's interrupts (node 1)},3/gray!12/solid/{housekeeping: kernel, logger, the rest}} {
  \draw[draw=gray!70,fill=\f,\d] (6.1,-0.5*\i-0.1) rectangle ++(0.35,0.3);
  \node[anchor=west,font=\footnotesize] at (6.5,-0.5*\i+0.05) {\t};
}
\end{tikzpicture}

\omcaption{A tuned host for chapter~4's two-socket fixture: each column is a physical core with its two hardware threads.
The boot line \texttt{isolcpus=21-23,29-31 nohz\_full=21-23 rcu\_nocbs=21-23 irqaffinity=0-15} with the siblings taken
offline, the card's interrupt queues bound to CPUs 16--20; the build's audit accepts this host (the \texttt{tuned}
fixture).}\label{fig:ll:linux-tuning:host}
\end{omfigure}

\section{Interrupts and their affinity}

\begin{definition}[Interrupt affinity]\label{def:ll:linux-tuning:irq}
The \emph{interrupt affinity}\index{interrupt affinity} of an interrupt line is the set of CPUs allowed to handle it, read and
written in \path{/proc/irq/N/smp_affinity_list}; the boot parameter \texttt{irqaffinity} sets the default for lines that no one
configures.
\end{definition}

An interrupt runs its handler on whichever CPU receives it, in the middle of whatever that CPU was doing: the network card
announcing a packet, a disk finishing a write, another CPU asking this one to flush its translation buffer. On an untuned
machine every line may go to every CPU: the audit of this laptop in the build below finds the default mask and all of its
active lines allowed on the CPUs that chapter~4 would give the hot path. The kernel documentation's instruction is short:
``isolate the IRQs whenever possible, so that they don't fire on the target CPUs''. Two details make it harder than writing a
mask. Many distributions run \texttt{irqbalance}, a daemon that periodically spreads interrupts across CPUs and would undo a
hand-made placement; the Red Hat real-time tuning guide says it is not needed where applications are bound to CPUs, and
disables it. And some interrupts belong on a hot core's socket but not on the core: the network card's receive queues should
interrupt a housekeeping CPU of the card's node (\cref{fig:ll:linux-tuning:host}), so that the kernel's packet processing
leaves the data in the socket's shared cache, where the hot threads read it, rather than across the interconnect
(chapter~4).

\begin{dated}{2026-09}{Distribution tuning profiles}\label{dat:ll:linux-tuning:rhel}
Red Hat's guide to optimising RHEL~9 for Real Time for low-latency operation (consulted September 2026) walks through the
same steps as this chapter: isolating CPUs with the TuneD real-time profile's \texttt{isolated\_cores} variable, disabling
\texttt{irqbalance} and binding interrupts by hand, locking memory with \texttt{mlock} and \texttt{mlockall}, and the
real-time throttling parameters of \cref{sec:ll:linux-tuning:rt}. Profiles and their variable names change between releases;
the kernel parameters they set are the stable part.
\end{dated}

\section{Tickless kernels and housekeeping}

\begin{definition}[Tickless kernel]\label{def:ll:linux-tuning:nohz}
A \emph{tickless kernel}\index{tickless kernel} (adaptive ticks, \texttt{nohz\_full}) stops the periodic scheduling-clock
interrupt on listed CPUs while each runs a single task, and moves the work that interrupt did for them (timekeeping, a residual
once-a-second tick, the callbacks of the read-copy-update mechanism) to housekeeping CPUs.
\end{definition}

The scheduling-clock interrupt, the tick, lets the kernel account time, run timers and decide whether to switch tasks,
100 to 1\,000 times a second depending on the build. A CPU with a single runnable task has nothing to switch to, and the
kernel's documentation on adaptive ticks puts the gain precisely: real-time applications ``improve their worst-case response
times by the maximum duration of a scheduling-clock interrupt''. The same document lists the conditions: the kernel must be
built with \texttt{CONFIG\_NO\_HZ\_FULL} (this laptop's is not), the CPU must run one task only, the thread must stay out of
the kernel (system calls and \omterm{def:ll:cpp-for-latency-i-memory:fault}{page faults} bring the noise back), and at least one CPU must keep ticking for timekeeping.
Read-copy-update callbacks, deferred work the kernel queues on each CPU, go with the tick: \texttt{nohz\_full} offloads them
from the listed CPUs, and \texttt{rcu\_nocbs} does so explicitly. The price is paid elsewhere: entering and leaving the kernel
costs more on a tickless CPU, and the housekeeping CPUs take the offloaded work.

What does a gap cost a message? The hook's thread spends its time either running or stopped, and a message arriving at a
random moment waits only if it arrives during a stop.

\begin{proposition}[The wait added by gaps]\label{prop:ll:linux-tuning:wait}
A thread that must handle each message immediately is stopped during gaps of lengths $g_1,\dots,g_n$ in an observation of
length $T$. A message arriving at a uniformly random moment of $[0,T]$ waits, beyond its own processing,
\[
\E[W] = \frac{1}{2T}\sum_{i=1}^{n} g_i^2, \qquad \P(W > x) = \frac{1}{T}\sum_{i=1}^{n} (g_i - x)^+ ,
\]
and the fraction of time during which it would wait at all is $\sum_i g_i/T$.
\end{proposition}

\begin{proof}
The message falls inside gap $i$ with probability $g_i/T$, and then its position in the gap is uniform, so it waits a
remaining time $U_i$ uniform on $[0, g_i]$, with mean $g_i/2$ and $\P(U_i > x) = (1 - x/g_i)^+$. Outside the gaps it does
not wait. Summing over the gaps gives $\E[W] = \sum_i (g_i/T)(g_i/2)$ and $\P(W > x) = \sum_i (g_i/T)(1 - x/g_i)^+$.
\end{proof}

The squares matter: the wait is dominated by the long gaps, not the frequent ones. A tick of \qty{2}{\micro\second} at
250~Hz steals 0.05\% of the time and adds half a nanosecond on average; a single gap of \qty{1}{\milli\second} a second steals
0.1\% but adds \qty{500}{\nano\second} on average and makes one message in two thousand wait more than half a millisecond.
On the laptop, pinned alone, the meter's gaps give an average added wait of about three microseconds (it moves by a factor
of two or more from run to run, with the run's longest gap: the squares at work), a wait beyond \qty{10}{\micro\second}
for about 0.45\% of messages and beyond \qty{100}{\micro\second} for 0.1 to 0.2\%: the host adds a 99.9th percentile of
over a hundred microseconds to a path whose own work takes a few microseconds or less (chapter~1). The proposition assumes messages arrive
independently of the gaps; bursts that coincide with the kernel's work (an interrupt from the same card that brings the
data) make it worse.

\section{Busy polling and kernel bypass}

\begin{definition}[Context switch, busy polling, kernel bypass]\label{def:ll:linux-tuning:polling}
A \emph{context switch}\index{context switch} is the kernel saving the state of the thread running on a CPU and restoring
another's; a thread that blocks waiting for data gives its CPU away and is switched back in when the data arrives. \emph{Busy
polling}\index{busy polling} checks for new data in a loop that never blocks, keeping the CPU. \emph{Kernel
bypass}\index{kernel bypass} moves the network card's receive and transmit queues into the application's address space, so
that packets reach the program without a system call, an interrupt or the kernel's network stack.
\end{definition}

A blocking receive is the default and the right choice for almost every program: the thread sleeps, the CPU does other work
or saves power, and the kernel wakes the thread when the packet has gone through the interrupt handler, the protocol stack and
the socket's buffer (\cref{fig:ll:linux-tuning:paths}, left). Each step costs microseconds, and the wake-up is the worst of
them: the sleeping thread's CPU may be idle in a low-power state, and the thread must be scheduled back in. \omterm{def:ll:linux-tuning:polling}{Busy polling}
removes the sleep: the thread asks for data in a loop, with a non-blocking call, and finds it as soon as the kernel has put it
in the socket (middle). The socket option \texttt{SO\_BUSY\_POLL} makes the kernel itself spin for up to a set number of
microseconds inside a blocking receive (socket(7)); the tutorial polls from user space instead, which works on any socket.
\omterm{def:ll:linux-tuning:polling}{Kernel bypass} removes the kernel: in the open-source DPDK, a poll-mode driver ``accesses the RX and TX descriptors directly
without any interrupts'', in a run-to-completion loop on a dedicated core (right). One Quant Book~14 covers the stacks and the
cards; what matters here is that each step to the right trades a CPU spinning at 100\% for microseconds.

\begin{omfigure}
\begin{tikzpicture}[font=\footnotesize,box/.style={draw,rounded corners=2pt,minimum width=2.9cm,minimum height=0.55cm,
  align=center,inner sep=2pt},arr/.style={-Stealth,thick}]
\foreach \x/\t in {0/{interrupts, blocking},4.4/{busy polling},8.8/{kernel bypass}} {\node[font=\footnotesize\bfseries] at (\x,0.7) {\t};}
\node[box,fill=gray!10] (n1) at (0,-4.4) {network card};
\node[box,fill=omThm!15] (i1) at (0,-3.3) {interrupt handler};
\node[box,fill=omThm!15] (s1) at (0,-2.2) {kernel network stack};
\node[box,fill=omThm!15] (b1) at (0,-1.1) {socket buffer, wake-up};
\node[box,fill=omDef!20] (a1) at (0,0) {thread asleep in \texttt{recv}};
\draw[arr] (n1) -- (i1); \draw[arr] (i1) -- (s1); \draw[arr] (s1) -- (b1);
\draw[arr,omThm] (b1) -- node[right,font=\scriptsize]{context switch} (a1);
\node[box,fill=gray!10] (n2) at (4.4,-4.4) {network card};
\node[box,fill=omThm!15] (s2) at (4.4,-3.0) {kernel network stack};
\node[box,fill=omThm!15] (b2) at (4.4,-1.6) {socket buffer};
\node[box,fill=omDef!20] (a2) at (4.4,0) {thread spinning};
\draw[arr] (n2) -- (s2); \draw[arr] (s2) -- (b2);
\draw[arr,omDef,dashed] (b2) -- node[right,font=\scriptsize,align=left]{non-blocking\\ call in a loop} (a2);
\node[box,fill=gray!10] (n3) at (8.8,-4.4) {network card};
\node[box,fill=gray!10] (r3) at (8.8,-2.2) {descriptor ring\\ in user memory};
\node[box,fill=omDef!20] (a3) at (8.8,0) {thread spinning};
\draw[arr] (n3) -- node[right,font=\scriptsize]{DMA} (r3);
\draw[arr,omDef,dashed] (r3) -- node[right,font=\scriptsize]{poll} (a3);
\end{tikzpicture}

\omcaption{Three ways for a packet to reach a thread. Left: the card interrupts a CPU, the kernel processes the packet, and
the sleeping thread is switched back in. Middle: the thread never sleeps and polls the socket. Right: the card writes into
queues mapped into the program, which polls them; the kernel is not on the path (DMA: \omterm{def:ll:multi-socket-machines-and-interconnects:pcie}{direct memory access}, chapter~4).
Shaded red: the kernel's work on the path.}\label{fig:ll:linux-tuning:paths}
\end{omfigure}

\Cref{fig:ll:linux-tuning:receive} measures the first two on the laptop's loopback interface, where no card is involved and the
difference is the sleep alone. A round trip of 64 bytes between two pinned threads takes about \qty{30}{\micro\second} at the
median with blocking receives and about \qty{2.3}{\micro\second} busy-polled, over ten times less. Two threads bouncing a byte
through a pair of pipes on the \emph{same} CPU complete a round trip in about \qty{1.5}{\micro\second}: two \omterm{def:ll:linux-tuning:polling}{context switches}
with their system calls, each well under a microsecond when the CPU is awake. The same ping-pong between two CPUs takes about
\qty{30}{\micro\second}: the cost is not the switch but waking a CPU that went idle while waiting, which is also why the
blocking receive is slow. \omterm{def:ll:linux-tuning:polling}{Kernel bypass}, which removes the system calls as well, needs a card that supports it, which this
laptop does not have.

\begin{omfigure}
\begin{tikzpicture}
\begin{axis}[width=12cm,height=5.8cm,ybar,bar width=7pt,ymode=log,log origin y=infty,ymin=500,ymax=1e7,
  symbolic x coords={udp-block,udp-spin,pipe-same,pipe-two},
  xtick=data,xticklabels={{UDP loopback,\\ blocking},{UDP loopback,\\ busy-polled},{pipe ping-pong,\\ one CPU},{pipe ping-pong,\\ two CPUs}},
  xticklabel style={align=center,font=\footnotesize},enlarge x limits=0.15,
  ylabel={ns per round trip (log scale)},tick label style={font=\small},label style={font=\small},grid=major,
  grid style={gray!25},legend columns=4,legend style={font=\footnotesize,at={(0.5,-0.33)},anchor=north,draw=none,fill=none,
  /tikz/every even column/.append style={column sep=6pt}},table/col sep=comma]
\addplot[fill=omDef!70,draw=omDef] table[x=key,y=p50]{figdata/low-latency/13-linux-tuning/measured_receive.csv};
\addplot[fill=omDef!35,draw=omDef] table[x=key,y=p99]{figdata/low-latency/13-linux-tuning/measured_receive.csv};
\addplot[fill=omThm!45,draw=omThm] table[x=key,y=p999]{figdata/low-latency/13-linux-tuning/measured_receive.csv};
\addplot[fill=gray!40,draw=gray] table[x=key,y=max]{figdata/low-latency/13-linux-tuning/measured_receive.csv};
\legend{p50,p99,p99.9,maximum}
\end{axis}
\end{tikzpicture}

\omcaption{Round trips of 64 bytes between two pinned threads over loopback UDP, receiving with a blocking call or with
non-blocking calls in a loop, and of one byte through two pipes with both threads on one CPU or on two; 100\,000 round trips
each. Measured on a laptop (Intel Core Ultra 7 155H) under WSL2, no isolated cores. Data:
\texttt{bench\_tuning.py}.}\label{fig:ll:linux-tuning:receive}
\end{omfigure}

The tails tell the rest: \omterm{def:ll:linux-tuning:polling}{busy polling} cuts the median, not the p99.9, which on this laptop is still tens of microseconds,
because the spinning thread meets the gaps of \cref{fig:ll:linux-tuning:hiccups} like any other. Polling and isolation go
together: a thread that polls on a CPU where the kernel still interrupts it is fast on average and slow in the tail, and an
isolated CPU whose thread sleeps in a blocking call wakes up slowly. The \omterm{def:ll:cpu-microarchitecture:dvfs}{idle states} are the third part: a CPU that is never
idle never enters a deep one, but the housekeeping CPUs that wake the hot path's helpers can, and the kernel documentation
lists \texttt{intel\_idle.max\_cstate=} as the way to limit their depth, and \texttt{idle=poll} as a stronger one that ``can
cause your CPU to overheat'' and disables turbo frequencies.

\section{Real-time priorities and memory locking}\label{sec:ll:linux-tuning:rt}

\begin{definition}[Real-time scheduling policy, memory locking]\label{def:ll:linux-tuning:rt}
A \emph{real-time scheduling policy}\index{real-time scheduling policy} (\texttt{SCHED\_FIFO} or \texttt{SCHED\_RR} in
Linux) gives a thread a static priority above every ordinary thread: a runnable real-time thread preempts any ordinary one and,
under \texttt{SCHED\_FIFO}, runs without time slicing until it blocks, yields or is preempted by a higher priority.
\emph{Memory locking}\index{memory locking} (\texttt{mlock}, \texttt{mlockall}) keeps a process's pages in physical memory:
they are never paged out, and with \texttt{MCL\_FUTURE} every page of every later mapping is made present when the mapping is
created.
\end{definition}

A real-time policy answers the shared-CPU curve of \cref{fig:ll:linux-tuning:hiccups}: the rival no longer gets a slice.
On an isolated core with a single thread there is no rival to preempt, and the policy brings a risk instead. A real-time
thread that spins forever would starve the CPU's kernel threads, so Linux throttles real-time threads by default: they may run
\qty{950}{\milli\second} of every second, and the remaining \qty{50}{\milli\second} go to ordinary tasks (the kernel's
documentation of real-time group scheduling). A spinning \texttt{SCHED\_FIFO} thread on a CPU that also has an ordinary task
runnable is therefore stopped for up to \qty{50}{\milli\second} once a second, a hiccup a thousand times longer than any on the
untuned laptop. Throttling can be disabled (\texttt{sched\_rt\_runtime\_us} set to $-1$), which the Red Hat guide calls
adequate only when the real-time tasks are well engineered. An ordinary process may not ask for the policy at all: on this
laptop the request fails with \texttt{EPERM}. The firm's design therefore spins ordinary threads on isolated cores, where
there is nothing to preempt, and keeps real-time priorities for threads that must preempt others on shared CPUs.

\omterm{def:ll:cpp-for-latency-i-memory:fault}{Page faults} are the synchronous noise of the kernel documentation, and chapter~6 measured them: the first write to a page
costs a microsecond or more. The mlockall manual page names the reason to lock: ``paging is one major cause of unexpected
program execution delays'' for real-time algorithms. On the laptop, writing one byte to each page of a fresh 32~MiB mapping
takes 8\,192 \omterm{def:ll:cpp-for-latency-i-memory:fault}{page faults} and about one to two microseconds a page. After \texttt{mlockall(MCL\_CURRENT | MCL\_FUTURE)}, the
same writes take no fault and about \qty{35}{\nano\second} a page: the cost has moved to the \texttt{mmap} call, which now
populates every page up front, at start-up, where it belongs. Locking is limited by \texttt{RLIMIT\_MEMLOCK}, 64~MiB for an
ordinary user on this machine: a process whose book and rings need a gigabyte must be given the limit (the build checks it),
or \texttt{mlockall} fails with \texttt{ENOMEM}.

\section{Tutorial: measure an untuned machine}\label{tut:ll:linux-tuning:tut}

\begin{tutorial}
\textbf{Goal.} Measure what this machine does to a thread that wants its CPU, and audit it against a placement plan.
\textbf{End state:} \cref{fig:ll:linux-tuning:hiccups,fig:ll:linux-tuning:receive}, the page-fault counts with and without
locking, and the audit's report for this machine next to the two fixtures'.
\begin{tutsteps}
\item \textbf{The hiccup meter.} A loop reads the \omterm{def:ll:cpu-microarchitecture:tsc}{time-stamp counter} and records every gap above a threshold in a buffer
reserved before the loop; it allocates nothing and makes no system call, so every gap it sees was imposed on it.
\omcode{code/low-latency/13-linux-tuning/cpp/ll_tuning.hpp}{25}{44}{The hiccup meter's loop, generic over its clock so that
the test can script one.}\label{lst:ll:linux-tuning:meter}
\item \textbf{Run it} with \texttt{python bench\_tuning.py}: five seconds pinned alone with \texttt{taskset}, five seconds
unpinned, and five seconds pinned beside a second spinning thread; the script converts the gaps into
\cref{fig:ll:linux-tuning:hiccups} and applies \cref{prop:ll:linux-tuning:wait} to them.
\item \textbf{Poll instead of sleeping.} The same receive, blocking or in a non-blocking loop:
\omcode{code/low-latency/13-linux-tuning/cpp/ll_tuning_bench.cpp}{109}{117}{A receive that sleeps and one that
spins.}\label{lst:ll:linux-tuning:receive}
\item \textbf{Lock the memory} and count the faults. The C++ test maps 8~MiB, touches it (2\,048 faults), then calls
\texttt{mlockall} and maps and touches another 8~MiB (none); if the process may not lock, it says so and skips that half.
\omcode{code/low-latency/13-linux-tuning/cpp/ll_tuning.hpp}{62}{71}{Counting page faults, and locking every present and future
page.}\label{lst:ll:linux-tuning:lock}
\item \textbf{Ask for a real-time policy} (\texttt{ll\_tuning\_bench privileges}): an ordinary process gets \texttt{EPERM}.
\item \textbf{Audit the machine}: \texttt{python firm\_tuneaudit.py /} plans the hot threads with \texttt{firm.coreplan} and
checks the host against the plan (next section).
\end{tutsteps}
\textbf{What to change next.} Run the meter on a CPU while another program (a compile, a test suite) runs on its hyper-threading
sibling: the gaps do not change, but count the loop's passes a second. On a Linux host you administer, boot with
\texttt{isolcpus} and \texttt{nohz\_full} on two CPUs, move their interrupts away, and run the meter there.
\end{tutorial}

\section{Build: the host audit}\label{bld:ll:linux-tuning:tuneaudit}

\texttt{firm.tuneaudit} reads the same files an engineer would, from \texttt{/proc} and \texttt{/sys} or from a fixture tree
with their layout, and checks each rule of this chapter against the plan of \texttt{firm.coreplan}: the hot CPUs isolated,
their siblings offline or isolated, the hot CPUs tickless and their read-copy-update callbacks offloaded, no interrupt line
(nor the default mask) allowed on a hot CPU, \texttt{irqbalance} not running, the \texttt{performance} frequency governor on
the hot CPUs, the \omterm{def:ll:cpu-microarchitecture:dvfs}{idle states} capped on the boot line, transparent \omterm{def:ll:memory-hierarchy-and-caches:tlb}{huge pages} only where asked for (\texttt{madvise}, with
no stalling \texttt{always} defragmentation, which the kernel documentation says makes an allocation ``stall on allocation
failure and directly reclaim pages and compact memory''), a locked-memory limit large enough, and real-time throttling off
if the hot threads are real-time.

\omcode{code/firm/tuneaudit/firm_tuneaudit.py}{72}{100}{The isolation, tick and interrupt rules of the
audit.}\label{lst:ll:linux-tuning:audit}

A setting the host does not expose is reported \texttt{unknown}, not passed: this laptop has no frequency governor to read, as
virtual machines often do not. Its report fails seven of the eleven rules (the hot CPUs neither isolated nor tickless, their
siblings busy, every interrupt allowed on them, no idle-state cap, a 64~MiB locked-memory limit), passes three and cannot judge
one. The \texttt{tuned} fixture, which describes \cref{fig:ll:linux-tuning:host}, passes all eleven; the \texttt{mistuned}
fixture fails nine, one per usual mistake.

\begin{build}
\bfield{Purpose} The check that a host is fit to run the hot path before a process starts on it: the deployment of chapter~25
runs it and refuses a host that fails, and its report goes with every latency measurement, so that a change of host tuning is
visible when the numbers move.
\bfield{Interface} Python \texttt{firm.tuneaudit}: \texttt{audit(root, plan, hot, memlock\_bytes, fifo)} returns a list of
\texttt{Finding(rule, status, detail)} with status \texttt{ok}, \texttt{fail} or \texttt{unknown}, one per rule of
\texttt{RULES} in order; \texttt{failures(findings)}; \texttt{report(findings)}; the command line audits a root and an interface
with \texttt{firm.coreplan}'s plan and exits non-zero on any failure.
\bfield{Rules} Read-only: the audit changes nothing on the host. Every rule is derived from the plan, not from a fixed CPU list;
an unreadable setting is \texttt{unknown}, never \texttt{ok}; details name the offending CPUs or interrupt lines.
\bfield{Acceptance tests} \texttt{code/firm/tuneaudit/}: the \texttt{tuned} fixture passes every rule and the
\texttt{mistuned} fixture fails exactly the nine expected ones, with the offending siblings, interrupt lines and governor named;
real-time threads with default throttling fail; a larger locking requirement fails only the mistuned limit; an empty tree reports
\texttt{unknown} where nothing can be read. The fixtures are written by \texttt{make\_tuneaudit\_fixture.py} for
\texttt{firm.coreplan}'s two-socket fixture.
\bfield{Stretch} Check that the card's interrupt queues are on housekeeping CPUs of the card's node; check the kernel's own
configuration (\texttt{CONFIG\_NO\_HZ\_FULL}) when it is exposed; emit the boot line and interrupt masks that would fix a
failing host.
\end{build}

\begin{omsources}
\item Linux kernel documentation: ``The kernel's command-line parameters'' (isolcpus, nohz\_full, rcu\_nocbs,
irqaffinity); ``CPU Isolation''; ``NO\_HZ: Reducing Scheduling-Clock Ticks''; ``Real-Time group scheduling'';
``Transparent Hugepage Support''.
\item Linux man pages \texttt{sched(7)}, \texttt{socket(7)} (\texttt{SO\_BUSY\_POLL}), \texttt{mlockall(2)}.
\item Red Hat, \emph{Optimizing RHEL 9 for Real Time for low latency operation}.
\item DPDK Programmer's Guide, ``Poll Mode Driver'' (release 24.07).
\end{omsources}

\section{Exercises}

\begin{exercise}[$\star$]\label{exo:ll:linux-tuning:1}
A kernel ticks 1\,000 times a second and each tick takes \qty{2}{\micro\second} on the CPU it interrupts. What fraction of a
spinning thread's time does it take, what average wait does it add to a message arriving at random
(\cref{prop:ll:linux-tuning:wait}), and what does \texttt{nohz\_full} remove from the worst case?
\end{exercise}

\begin{exercise}[$\star$]\label{exo:ll:linux-tuning:2}
A process whose rings and books occupy 3~GiB calls \texttt{mlockall(MCL\_CURRENT | MCL\_FUTURE)} as an ordinary user on this
laptop. What happens, and what must change?
\end{exercise}

\begin{exercise}[$\star$]\label{exo:ll:linux-tuning:3}
A host has 32 cores with two hardware threads each; the sibling of CPU $c < 32$ is $c + 32$. The plan puts four hot threads on
CPUs 12--15. Write the \texttt{isolcpus} and \texttt{nohz\_full} lists, and say what to do with the siblings.
\end{exercise}

\begin{exercise}[$\star\star$]\label{exo:ll:linux-tuning:4}
A thread loses 1\,000 gaps a second of \qty{10}{\micro\second} each. What fraction of its time is stolen, what is the mean
added wait, and what fraction of messages wait more than \qty{5}{\micro\second}? Compare with one gap a second of
\qty{10}{\milli\second}, which steals the same time.
\end{exercise}

\begin{exercise}[$\star\star$]\label{exo:ll:linux-tuning:5}
\omterm{def:ll:linux-tuning:polling}{Busy polling} keeps a core at 100\% whether or not data arrives. When is that a price worth paying, and what does it cost
besides the core?
\end{exercise}

\begin{exercise}[$\star\star$]\label{exo:ll:linux-tuning:6}
A team gives its spinning feed thread \texttt{SCHED\_FIFO} priority on a CPU where a monitoring agent also runs, with the
default throttling. What gap does the feed thread suffer, how often, and what share of its time? What are the two fixes?
\end{exercise}

\begin{exercise}[$\star\star\star$]\label{exo:ll:linux-tuning:7}
\emph{Coding.} Add a rule \texttt{nic\_irqs} to \texttt{firm.tuneaudit}: the interrupt lines of the plan's network card (read
their names from a fixture \path{/proc/interrupts}) must be allowed only on housekeeping CPUs of the card's node. Extend both
fixtures and the tests.
\end{exercise}

\begin{exercise}[$\star\star\star$]\label{exo:ll:linux-tuning:8}
\emph{Find the flaw.} ``We isolated CPUs 2--5 with \texttt{isolcpus} and started the feed handler with \texttt{taskset -c 2-5}
so that its four threads would share them. Its latency got worse.''
\end{exercise}

\section{Problem: The Hiccup}

\begin{problem}[{Weekend problem --- what the host takes from a spinning thread}]\label{pb:ll:linux-tuning:1}
The hiccup meter ran for five seconds on this laptop in three situations
(\texttt{measured\_hiccup\_summary.csv}). A tuned host is modelled with stated assumptions, not measurements: with the tick
running, 250 interruptions a second of \qty{2}{\micro\second}; with the hot CPU isolated and tickless, one interruption of
\qty{2}{\micro\second} a second.

\medskip\noindent\textbf{Part I --- The meter.}
\begin{enumerate}
\item Why must the meter's loop never allocate memory or call the kernel?
\item From the summary, how long does one pass of the loop take, and why is the \qty{200}{\nano\second} threshold safe?
\item What fraction of its time does the thread lose pinned alone, and what is its worst gap?
\item What fraction does it lose sharing its CPU with a second spinner, and why about that?
\end{enumerate}

\medskip\noindent\textbf{Part II --- What a message sees.}
\begin{enumerate}[resume]
\item Using \cref{prop:ll:linux-tuning:wait}, what does the summary give for the mean added wait pinned alone?
\item What fraction of messages wait more than \qty{10}{\micro\second}, and more than \qty{100}{\micro\second}?
\item In the shared case, what is the median added wait of a message?
\item Why do long gaps matter more than frequent ones?
\end{enumerate}

\medskip\noindent\textbf{Part III --- The tuned host.}
\begin{enumerate}[resume]
\item With the 250~Hz tick of \qty{2}{\micro\second}: stolen fraction and mean added wait.
\item With the tick: what fraction of messages waits more than \qty{1}{\micro\second}?
\item Isolated and tickless: stolen fraction and mean added wait.
\item Where does the residual once-a-second tick of a tickless CPU run?
\item If the feed thread were \texttt{SCHED\_FIFO} with default throttling and a runnable ordinary task on its CPU, what
would the worst gap and stolen fraction be?
\item Why can the laptop not be tuned to the model?
\end{enumerate}

\medskip\noindent\textbf{Part IV --- The verdict.}
\begin{enumerate}[resume]
\item State the \emph{named result}: the fraction of wall time a spinning thread is not running, and its worst gap, on this
laptop against the tuned host.
\item Which of the audit's rules does this laptop fail (\texttt{measured\_audit.csv})?
\item How would you run the meter on a production host without disturbing trading?
\item What does the meter not tell you about a gap, and how would you find it?
\item What should be recorded with every latency measurement, following this chapter?
\item In one sentence: what does tuning buy, and what does it cost?
\end{enumerate}
\end{problem}

\section{Interview questions}

\begin{interviewq}[$\star$ \iqroles{developer}]\label{iq:ll:linux-tuning:1}
What do \texttt{isolcpus}, \texttt{nohz\_full} and \texttt{rcu\_nocbs} do, and why are they used together?
\end{interviewq}

\begin{interviewq}[$\star\star$ \iqroles{developer}]\label{iq:ll:linux-tuning:2}
Your pinned, spinning thread shows a latency spike every few milliseconds. How do you find the cause?
\end{interviewq}

\begin{interviewq}[$\star\star$ \iqroles{developer}]\label{iq:ll:linux-tuning:3}
Blocking receive, \omterm{def:ll:linux-tuning:polling}{busy polling}, or \omterm{def:ll:linux-tuning:polling}{kernel bypass}: what does each cost, and when would you choose each?
\end{interviewq}

\begin{interviewq}[$\star\star$ \iqroles{developer}]\label{iq:ll:linux-tuning:4}
Should the hot threads of a trading process run under \texttt{SCHED\_FIFO}? What can go wrong?
\end{interviewq}

\begin{interviewq}[$\star\star$ \iqroles{developer}]\label{iq:ll:linux-tuning:5}
Where should the network card's interrupts go on a tuned host, and why not simply away from every CPU of its socket?
\end{interviewq}

\begin{interviewq}[$\star\star\star$ \iqroles{developer, researcher}]\label{iq:ll:linux-tuning:6}
A colleague reports that tuning halved the median tick-to-trade but did nothing to the p99.9. What would you look at?
\end{interviewq}
